\documentclass{article}
\usepackage[top=2.3cm, bottom=2.3cm, left=2.1cm, right=2.1cm]{geometry}
\usepackage{graphicx} % Required for inserting images
\usepackage{threeparttable}
\usepackage{booktabs}
\usepackage{rotating}
\usepackage{multirow}
\usepackage{hyperref}
\usepackage{float}
\usepackage{amsmath}
\usepackage{mathtools}
\usepackage{amssymb}
\usepackage{dsfont}
\usepackage{booktabs}
\usepackage{amsthm}
\usepackage{graphicx}
\usepackage[dvipsnames]{xcolor}
\usepackage{setspace}
\usepackage{babel}
\usepackage{pstricks}
\usepackage{pstricks-add}
\let\clipbox\relax
\usepackage{tikz}
%\usepackage{subfig}
\usepackage{adjustbox}
\usepackage{enumitem}
\usepackage{bbm}
\usepackage{hyperref}
\usepackage{natbib}
\usepackage{multirow}
%\usepackage{pgfplots}
\usepackage{accents}
\usepackage{caption}
\usepackage{subcaption}
\usepackage{lscape}
\hypersetup{colorlinks,linkcolor={blue},citecolor={blue},urlcolor={red}}
%\usepackage[capposition=top]{floatrow}
\usepackage{microtype}
\usepackage[default,regular,black]{sourceserifpro}
\usepackage[T1]{fontenc}

\newcommand{\Align}{\mathrm{Align}}
\newcommand{\FA}{\mathrm{FA}}
\newcommand{\EC}{\mathrm{EC}}

\begin{document}

\title{Testing Industrial Policymakers: Evidence from Brazil's BNDES}
\author{}
\date{\today}

\maketitle

\begin{abstract}
    Industrial policy. Everyone is doing it. Is it effective?
\end{abstract}

\normalsize

\newpage
\clearpage
\setcounter{page}{1}
\setlength{\parskip}{10pt}
\setlength{\parindent}{10pt}

\section{Shift-Share Idea}

Our identification strategy exploits the fact that political alignment shocks differentially affect sectors based on their pre-existing exposure to party-affiliated firms. If the policymaker is optimally allocating resources across sectors, then exogenous perturbations to this allocation---induced by changes in which party controls the mayoralty, governorship, or presidency---should have no first-order effect on municipal output. We implement this test using a shift-share instrumental variables design, where the ``shift'' is the change in political alignment at the municipality level and the ``share'' is each sector's baseline exposure to the newly aligned party.

The identifying assumption is that conditional on sector, municipality and time fixed effects, the interaction of political alignment shocks with pre-determined sectoral party exposure is orthogonal to unobserved determinants of municipality output. This assumption would be violated if, for example, parties systematically gain power in municipalities where their affiliated sectors were about to experience productivity growth for reasons unrelated to BNDES credit. We address this concern in several ways. First, we show that our instruments predict BNDES sectoral credit but not other municipal revenue sources or federal transfers. Second, we conduct placebo tests using leads of political alignment, which should have no predictive power if the parallel trends assumption holds. Third, we implement the Borusyak et al. (2022) shock-level inference procedure, which accounts for potential correlation in shocks across municipalities.

When party $p$ becomes aligned in a place–time cell $(m,t)$, BNDES tilts relatively more credit toward firms connected to $p$. The sectors of municipality $m$ that (before the elections) were home to more $p$-connected companies obtain a greater boost, plausibly exogenous within the municipality, to their share $s_{jmt}$. That is the “shift” (alignment) times “share” (baseline party exposure).

\section{Specifications}


\subsection{Firm-level, Levels}

The sector-level instruments $Z^{\ell}_{jmt}$ aggregate firm-level political connections within each sector. To validate that the political channel operates at the micro level, we estimate firm-level first-stage specifications that test whether firms whose owners are aligned with officeholders receive more BNDES credit.

\paragraph{Baseline party exposure.}
Given $L_{f,p,t}$ the number of owners of firm $f$ affiliated with party $p$ at time $t$,
and $L_{f,t}$ the total number of owners of firm $f$ at time $t$
(including owners not affiliated with any party),
we define the baseline exposure of firm $f$ to party $p$ as
\[
  \omega_{fp,\tau} \equiv \frac{L_{f,p,\tau}}{L_{f,\tau}},
  \qquad
  \sum_{p} \omega_{fp,\tau} \leq 1,
\]
where $\tau = \tau(t)$ denotes the pre-election baseline year for the electoral cycle
containing $t$.%
\footnote{Cycle-specific baselines: for mayors elected in year $e$,
$\tau = e - 1$ (pre-election year);
for governors and presidents, analogously.
As a robustness exercise we fix $\tau = 2002$, the earliest year available in the data, for all cycles.}
The residual $1 - \sum_p \omega_{fp,\tau}$ represents the share of owners
not affiliated with any political party;
these owners contribute zero to the instrument because they have no alignment shock.

\paragraph{Political shocks.}
For office $\ell\in\{\mathrm{mayor},\mathrm{governor},\mathrm{president}\}$, we define the political shock $\Align^{\ell}_{mpt}$ as an indicator of whether party $p$ holds office $\ell$ in the jurisdiction to which municipality $m$ belongs at period $t$. Throughout, $\ell$ indexes offices; alignment can be defined at the party level or at the coalition level (i.e., whether $p$ belongs to the governing coalition for office $\ell$). We present both variants in the empirical results. We assume that alignment shocks last the full term length (four years): if a party becomes aligned at inauguration year $\tau+1$, then $\Align^{\ell}_{mpt} = 1$ for $t \in \{\tau+1, \tau + 2, \tau + 3, \tau + 4\}$.

\paragraph{Firm-level shift-share instrument.}
Define the firm-level levels instrument as
\[
  \FA^{\ell}_{fmt} \equiv \sum_{p} \omega_{fp,\tau} \cdot \Align^{\ell}_{mpt}.
\]
This is the share of firm $f$'s owners who are affiliated with the party currently
holding office $\ell$ in the jurisdiction to which municipality $m$ belongs.
Because $\omega_{fp,\tau}$ does not vary across municipalities,
variation in $\FA^{\ell}_{fmt}$ across firms in the same $(m,t)$ cell
is driven purely by differences in owner-party composition.

\paragraph{Extensive margin (levels).}
We first ask whether alignment affects the probability of receiving any BNDES loan.
Using the full firm panel, we estimate the linear probability model
\[
  \mathbf{1}(\text{BNDES}_{fmt} > 0) = \sum_{\ell} \lambda^{\text{ext}}_{\ell}\, \FA^{\ell}_{fmt}
  + \gamma_{f} + \alpha_{mt} + u_{fmt},
\]
where $\gamma_f$ are firm fixed effects
and $\alpha_{mt}$ are municipality$\times$year fixed effects.
Standard errors are two-way clustered by firm and municipality.
The primary specification uses analytic regression weights equal to firm employment
($n_{\text{employees},ft}$), so that $\lambda^{\text{ext}}_\ell$ estimates the
employment-weighted average effect of alignment through office $\ell$
on the probability of receiving BNDES credit.
An unweighted specification serves as robustness.

\paragraph{Intensive margin (levels).}
Conditional on receiving a positive loan in a given year, we estimate
\[
  \log(\text{BNDES}_{fmt}) = \sum_{\ell} \lambda^{\text{int}}_{\ell}\, \FA^{\ell}_{fmt}
  + \gamma_{f} + \alpha_{mt} + u_{fmt},
  \qquad \text{BNDES}_{fmt} > 0,
\]
where $\lambda^{\text{int}}_\ell$ captures the semi-elasticity of loan size
with respect to alignment, among firms that receive any credit.
Fixed effects, clustering, and employment weighting follow the extensive-margin specification.
Comparing $\lambda^{\text{ext}}_\ell$ and $\lambda^{\text{int}}_\ell$ reveals
whether alignment operates primarily through access to credit or through the size of loans granted.

\subsection{Firm-level, Changes}

\paragraph{Firm-level alignment turnover instrument.}
To exploit only the variation induced by political turnover, define
\[
  \Delta\FA^{\ell}_{fmt} \equiv \sum_{p} \omega_{fp,\tau} \cdot \Delta\Align^{\ell}_{mpt},
\]
where $\Delta\Align^{\ell}_{mpt} = \Align^{\ell}_{mpt} - \Align^{\ell}_{mp,t-1}$
is the alignment turnover shock.%
\footnote{\label{fn:not-first-diff}$\Delta\FA^{\ell}_{fmt}$ is constructed from alignment turnover,
not as the first difference of $\FA^{\ell}_{fmt}$.
The two coincide only when baseline weights $\omega_{fp,\tau}$ are constant across
adjacent years---which holds within an electoral cycle but not at cycle boundaries.}

\paragraph{Extensive margin (changes).}
We first estimate whether political turnover affects the probability of receiving a loan:
\[
  \Delta\mathbf{1}(\text{BNDES}_{fmt} > 0) = \sum_{\ell} \lambda^{\text{ext}}_{\ell}\,\Delta\FA^{\ell}_{fmt}
  + \gamma_{f} + \alpha_{mt} + u_{fmt},
\]
where $\Delta\mathbf{1}(\text{BNDES}_{fmt} > 0)
= \mathbf{1}(\text{BNDES}_{fmt} > 0) - \mathbf{1}(\text{BNDES}_{fm,t-1} > 0)$.
This specification is more demanding than the levels version: identification relies exclusively on election-year political turnover rather than cross-sectional differences in alignment status.

\paragraph{Intensive margin (changes).}
Conditional on positive lending in both $t$ and $t-1$, we estimate
\[
  \Delta\log(\text{BNDES}_{fmt}) = \sum_{\ell} \lambda^{\text{int}}_{\ell}\,\Delta\FA^{\ell}_{fmt}
  + \gamma_{f} + \alpha_{mt} + u_{fmt},
  \qquad \text{BNDES}_{fmt} > 0,\; \text{BNDES}_{fm,t-1} > 0,
\]
where $\Delta\log(\text{BNDES}_{fmt}) = \log(\text{BNDES}_{fmt}) - \log(\text{BNDES}_{fm,t-1})$.
Fixed effects, clustering, and employment weighting follow the extensive-margin specification.

\subsection{Sector-level, Levels}

\paragraph{Sector-level baseline exposure.}
For sector $j$ in municipality $m$, define the sector-party exposure weight as
\[
  w_{jmp,\tau} \equiv
  \frac{\sum_{f\in\mathcal{F}(j,m)} L_{f,p,\tau}}
       {\sum_{f\in\mathcal{F}(j,m)} L_{f,\tau}},
\]
where $\mathcal{F}(j,m)$ is the set of firms in sector $j$ of municipality $m$.
This is the owner-count-weighted average of firm-level exposures:
$w_{jmp,\tau} = \sum_{f\in\mathcal{F}(j,m)} (L_{f,\tau}/N_{jm,\tau})\,\omega_{fp,\tau}$,
where $N_{jm,\tau} = \sum_{f\in\mathcal{F}(j,m)} L_{f,\tau}$.
As at the firm level, unaffiliated owners enter the denominator but contribute zero
to the instrument; $\sum_p w_{jmp,\tau} \leq 1$.

\paragraph{Levels instrument.}
The sector-level levels instrument is
\[
  Z^{\ell}_{jmt} \equiv \sum_{p} w_{jmp,\tau}\,\Align^{\ell}_{mpt}.
\]
Because $\Align^{\ell}_{mpt}$ is constant across sectors $j$ within $(m,t)$,
variation in $Z^{\ell}_{jmt}$ across sectors is driven by baseline exposure $w_{jmp,\tau}$.

\paragraph{First stage and exclusion.}
We estimate
\[
  s_{jmt} = \sum_{\ell} \lambda_{\ell}\, Z^{\ell}_{jmt}
  + \phi\,\EC_{jm,\tau} + \gamma_{jm} + \alpha_{jt} + u_{jmt},
\]
where $s_{jmt}$ is sector $j$'s share of BNDES credit in municipality $m$ at time $t$,
$\EC_{jm,\tau} \equiv \sum_p w_{jmp,\tau}$ is the sector-level exposure control
(measuring overall political connectedness of the sector-municipality cell),
$\gamma_{jm}$ are municipality$\times$sector fixed effects,
and $\alpha_{jt}$ are sector$\times$year fixed effects.
Standard errors are two-way clustered by municipality and sector.

\paragraph{Aggregation link.}
The sector-level instrument is the owner-count-weighted aggregation of firm-level instruments:
\[
  Z^{\ell}_{jmt}
  = \sum_{f\in\mathcal{F}(j,m)} \frac{L_{f,\tau}}{N_{jm,\tau}}\,\FA^{\ell}_{fmt}.
\]
The aggregation weight $L_{f,\tau}/N_{jm,\tau}$ is each firm's share of total owners
in the sector-municipality cell.
The employment-weighted counterpart---aggregating $\FA^{\ell}_{fmt}$ with weights
proportional to $n_{\text{employees},ft}$---produces a distinct object that corresponds
to the estimand of the employment-weighted firm first stage.
We verify this aggregation identity diagnostically by collapsing firm instruments
within each $(j,m,t)$ cell and confirming that the collapsed quantity satisfies
the instrument's support bounds.
Therefore, the firm-level F-statistics provide a micro-level validation of the sector-level instrument relevance: if political alignment does not predict firm-level lending, the sector-level instrument cannot have a first stage.

\subsection{Sector-level, Changes}

\paragraph{Shift-share changes instrument.}
The sector-level changes instrument is
\[
  \Delta Z^{\ell}_{jmt} \equiv \sum_{p} w_{jmp,\tau}\,\Delta\Align^{\ell}_{mpt}.
\]
As with the firm-level analog $\Delta\FA$,
this is constructed from alignment turnover, not as the first difference of $Z^{\ell}_{jmt}$.%
\footnote{See footnote~\ref{fn:not-first-diff} above.
Within an electoral cycle, $w_{jmp,\tau}$ is constant and the two coincide.}

\paragraph{First stage.}
We estimate
\[
  \Delta s_{jmt} = \sum_{\ell} \lambda_{\ell}\,\Delta Z^{\ell}_{jmt}
  + \phi\,\EC_{jm,\tau} + \gamma_{jm} + \alpha_{jt} + u_{jmt}.
\]
Because shares sum to unity within each municipality-year
($\sum_j s_{jmt} = 1$), changes satisfy the simplex constraint
$\sum_j \Delta s_{jmt} = 0$ in municipality-years with positive total BNDES
in both $t$ and $t-1$.
We therefore drop the sector with the largest mean share ($j_0$)
and interpret coefficients relative to it.

\subsection{Two Complementary Pipelines}

The firm-level and sector-level specifications estimate related but distinct objects.
The sector pipeline uses owner-count exposure weights $w_{jmp,\tau}$
and does not employ analytic regression weights.
The firm pipeline uses firm-specific exposure $\omega_{fp,\tau}$
with employment-weighted regression estimation ($n_{\text{employees},ft}$ as weights),
so that the firm-level estimand is the employment-weighted average treatment effect.

The employment-weighted aggregation of firm instruments within a sector-municipality cell
recovers an employment-weighted sector-level instrument:
\[
  \sum_{f\in\mathcal{F}(j,m)}
    \frac{n_{\text{employees},ft}}{\sum_{f'} n_{\text{employees},f't}}\,
    \FA^{\ell}_{fmt}
  = \sum_p w^{\text{emp}}_{jmp,\tau}\,\Align^{\ell}_{mpt},
\]
where $w^{\text{emp}}_{jmp,\tau}
= \sum_f n_{\text{employees},ft} \cdot \omega_{fp,\tau}
  / \sum_f n_{\text{employees},ft}$
is the employment-weighted exposure share.
This quantity differs from the owner-count instrument $Z^{\ell}_{jmt}$
whenever the employment distribution across firms differs from the owner-count distribution.
Both are valid shift-share instruments exploiting the same underlying political variation;
we present them as complementary views of the sector-level first stage.

\end{document}

